\renewcommand{\tCoverLine}{El arcade en MSX}
\renewcommand{\tComputerText}{\textbf{MSX1 o posterior}, con una ranura para el cartucho del juego y otra para la V9990.}
\renewcommand{\tMemoryText}{\textbf{64 KB de RAM recomendados}.}
\renewcommand{\tVideoText}{Se necesita una tarjeta \textbf{Yamaha V9990}. Los créditos de arranque se ven por la salida de vídeo del MSX. La ilustración, el título y el juego usan la salida de la V9990.}
\renewcommand{\tSoundText}{El PSG integrado del MSX. No hace falta un cartucho de sonido.}
\renewcommand{\tControlsText}{Un jugador. Teclado o joystick de dos botones en el puerto 1.}
\renewcommand{\tSetGame}{Inserta el \textbf{cartucho de \gameTitle{}} en otra ranura libre.}
\renewcommand{\tSetOn}{Enciende el MSX. Su salida de vídeo muestra los créditos y \texttt{V9990 found!}. Si indica \texttt{V9990 NOT found!}, apaga el equipo, comprueba la tarjeta y vuelve a probar.}
\renewcommand{\tBlackScreen}{\textbf{Si la pantalla del juego se queda en negro}, mira el aviso de arranque en la salida de vídeo del MSX y después cambia el monitor a la salida de la V9990. En openMSX, añade la extensión V9990 y carga la ROM como ASCII16.}
\renewcommand{\tCreditsTitle}{9. Créditos y derechos}
\renewcommand{\tSourceNote}{Código fuente disponible en GitHub. \textbf{¡Mejora el juego si te apetece!}}
\newcommand{\tCreditsExtra}{\item La portada y la etiqueta usan la ilustración de apertura del proyecto. Las capturas proceden de la versión para MSX.}
\newcommand{\tRelease}{Versión}
